% ============ Solutions and hints (Part 9: §25, §26) ============
% 原书 p.296; 原书对 §25 的 Ex.25.1, 25.3 未给解答, 编号按原书跳号.

\section*{\S\ref{sec:25}}

\noindent\textbf{25.2.}\quad If $xy = 0$ we prove that $y \in \bigcap_n x^nA$; suppose that $y \in x^nA$,
and set $y = x^nz$; then $0 = D(x^{n+1}z) = (n + 1)x^nz + x^{n+1}Dz$, and
$y \in x^{n+1}A$.

\noindent\textbf{25.4.}\quad $0 \to I \to A \otimes_k A \to A \to 0$ is a split exact sequence, so that
$0 \to I \otimes_k k' \to A \otimes_k A \otimes_k k' = A' \otimes_{k'} A' \to A' \to 0$ is also
exact, and hence $\Omega_{A'/k'} = (I/I^2) \otimes_k k' = \Omega_{A/k} \otimes_k k'$. For $A_S$, use the
fact that $A_S$ is 0-\'etale over $A$ and Theorem~\hyperref[thm:25.1]{1}.

\noindent\textbf{25.5.}\quad See Theorem~\hyperref[thm:27.3]{27.3}.

\section*{\S\ref{sec:26}}

\noindent\textbf{26.1.}\quad (i) Let $\alpha \in K \cap K'$; then $1, \alpha \in K$ are linearly dependent over
$K'$, hence also over $k$, so $\alpha \in k$. (ii) Assume that $u_1, \ldots, u_n \in K$ are
linearly independent over $k'$, and that $\sum \xi_iu_i = 0$ with
$\xi_i \in k'(K')$; we show that $\xi_i = 0$. Clearing denominators, we can assume
that $\xi_i \in k'[K']$. Choosing a basis $\{\omega_j\}$ of $K'$ over $k$ we can write
$\xi_i = \sum_j c_{ij}\omega_j$ with $c_{ij} \in k'$. Then since
$\sum_{i,j}c_{ij}\omega_ju_i = 0$ we get $\sum_i c_{ij}u_i = 0$, therefore
$c_{ij} = 0$ for all $i, j$.

\noindent\textbf{26.2.}\quad It is enough to show that $K((T))$ and $L^p((T^p))$ are linearly
disjoint over $K^p((T^p))$. Assume that $\omega_1(T), \ldots, \omega_r(T) \in K((T))$ are
linearly independent over $K^p((T^p))$, and that $\sum \varphi_i\omega_i = 0$ with
$\varphi_i \in L^p((T^p))$; we show that $\varphi_i = 0$ for all $i$. Clearing
denominators, we can assume that $\omega_i \in K[T]$ and
$\varphi_i \in L^p[T^p]$. Letting $\{\eta_\lambda\}$ be a basis of $L$ over $K$ we can in a
unique way write $\varphi_i = \sum_\lambda \varphi_{i\lambda}(T^p)$ with
$\varphi_{i\lambda}(T^p) \in K^p[T^p]$. Here $\sum_{\lambda,i}\varphi_{i\lambda}\omega_i$ is in
general an infinite sum, but only a finite number of terms appear in the sum
for the coefficient of some monomial in the $T$'s, so that the sum is
meaningful. Then $\sum_\lambda\big(\sum_i \varphi_{i\lambda}(T^p)\omega_i(T)\big) = 0$, so that
$\sum_i \varphi_{i\lambda}(T^p)\omega_i = 0$ for all $\lambda$, so
$\varphi_{i\lambda} = 0$ for all $i, \lambda$.
